\documentclass[11pt]{article}
\usepackage[margin=1.2in]{geometry}
\usepackage{amsmath,amssymb}
\title{A Kakeya Set Miniature: Tubes in the Unit Cube}
\author{Test Fixture}
\date{}
\begin{document}
\maketitle

A \emph{Kakeya set} in $\mathbb{R}^{3}$ contains a unit line segment in every
direction. We study the Besicovitch--Kakeya conjecture at finite scale: a
family $\mathbb{T}$ of $\delta$-tubes pointing in $\delta$-separated
directions satisfies the Wolff axioms, and we measure the volume of its union
on a voxel grid.

\section{Setup}
Take $\delta = 2^{-4}$ and $N = 96$ tubes through directions sampled on the
Fibonacci hemisphere. The Kakeya maximal estimate we probe numerically is
\[
\Bigl|\bigcup_{T\in\mathbb{T}} Y(T)\Bigr| \;\ge\; \delta^{\varepsilon}
\lambda^{K} \sum_{T\in\mathbb{T}} |T|.
\]
Dragging the packing parameter from separated to sticky, all tubes converge
without changing directions; the multiplicity of the central voxel grows
toward $N$.

\section{Observation}
At finite resolution the union volume decreases as the family packs toward
the Kakeya configuration, in line with the sticky Kakeya phenomenon.

\end{document}
